% =====================================================================
%  MATHEMATICAL TOOLKIT  (precedes the Measurement chapter)
% =====================================================================
\partdivider{A}{The Language and Kinematics}{Before any force or law, we need the vocabulary of motion, and the mathematics of direction. Everything mechanics says is written in these terms.}{\chref{ch:measurement} Measurement\\[1.2mm]
\chref{ch:straightline} Straight-Line Motion\\[1.2mm]
\chref{ch:vectors} Vectors\\[1.2mm]
\chref{ch:projectiles} Projectiles\\[1.2mm]
\chref{ch:relative} Relative Motion\\[1.2mm]
\chref{ch:circkin} Circular Kinematics}

\ncertchapter{0}{Mathematical Toolkit}\label{ch:toolkit}

Physics is stated in mathematics, and a derivation that keeps stopping to justify its algebra is
hard to follow. This chapter collects, in one place, the mathematical tools the rest of the volume
relies on, so that later arguments need not break off to explain a technique. Nothing here is examined
for its own sake. What is examined, constantly, is whether you can reach for the right tool at the
right moment, and that is a skill best built before the physics arrives to distract you.

There are twelve tools, and each one earns its place by appearing in a later chapter; each entry ends
with a pointer to where it is used. Read the chapter once now, quickly, for orientation. Then return
to it whenever a step in a derivation feels unfamiliar. The test of whether you have understood a
tool is not whether you can state it but whether you can say what it is \emph{for}: the dot product
answers ``how much along?'', the derivative answers ``how fast, right now?'', the integral answers
``how much, in total?''.

One honest warning. The calculus here is used, not proved. Where a result such as the derivative of
$\sin x$ is quoted, a mathematics course will supply the proof; this book supplies the physical
meaning and the habit of use. That division of labour is deliberate, and it is the standard one.

\section{Angles, Radians and the Trigonometric Ratios}

Two facts about angles are used on almost every page of this book: that calculus requires them to be
measured in radians, and that trigonometry is mostly used to split one vector into two perpendicular parts.
(A \emph{vector} is a quantity with a direction as well as a size, such as a displacement or a force;
its size is its \emph{magnitude}. \chref{ch:vectors} makes the idea precise.)

\begin{defbox}{Radians: the only angle measure that works in calculus}
An angle in \textbf{radians} is defined as the arc it subtends divided by the radius:
\[
\theta=\frac{\text{arc}}{\text{radius}}=\frac{s}{R}\qquad\Longrightarrow\qquad s=R\theta.
\]
Because it is a length divided by a length, a radian is \emph{dimensionless}. What makes it the only
measure fit for calculus is that, in radians and only in radians, $\sin\theta\approx\theta$ for small
$\theta$ and $\frac{d}{d\theta}\sin\theta=\cos\theta$; in degrees every such result carries an extra
factor $\pi/180$. A full turn is $2\pi$ radians, so
\[
180^\circ=\pi\ \text{rad},\qquad 1\ \text{rad}\approx57.3^\circ.
\]
\emph{Used in:} arc length and angular displacement (\chref{ch:circkin}), $v=\omega r$ (\chref{ch:circkin}), and every
rotational formula in \chsref{ch:rotation}{ch:rolling}. \textbf{An angle that stands on its own in a
formula} (as in $s=R\theta$ or $\omega=\theta/t$), rather than inside $\sin$, $\cos$ or $\tan$,
\textbf{must be in radians}; convert first. Inside a trigonometric function either unit is safe,
provided the calculator is set to match.
\end{defbox}

\begin{defbox}{The trigonometric ratios, and resolving a vector}
In a right triangle with hypotenuse of length $A$ and an angle $\theta$ (the sides named relative to $\theta$),
\[
\sin\theta=\frac{\text{opposite}}{\text{hyp}},\qquad
\cos\theta=\frac{\text{adjacent}}{\text{hyp}},\qquad
\tan\theta=\frac{\sin\theta}{\cos\theta},
\]
together with the identity that follows from Pythagoras,
\[
\boxed{\;\sin^{2}\theta+\cos^{2}\theta=1\;}
\]
\textbf{The one use that dominates this book} is \emph{resolution}: replacing a vector by two
perpendicular pieces that can be handled independently, the two legs of a right triangle whose hypotenuse
is the vector (Figure~\ref{fig:tkcomp}).
\[
A_x=A\cos\theta,\qquad A_y=A\sin\theta,\qquad
A=\sqrt{A_x^{2}+A_y^{2}},\qquad \tan\theta=\frac{A_y}{A_x},
\]
where $\theta$ is measured from the $x$-axis and $A_x$, $A_y$ are the \emph{components} along the $x$- and
$y$-axes. \emph{Used in:} projectile motion (\chref{ch:projectiles}), every inclined-plane problem
(Chapters~\ref{ch:newton}--\ref{ch:friction}), and the components
of weight on a slope, $mg\sin\theta$ and $mg\cos\theta$.
\end{defbox}

\begin{figure}[htbp]
\centering
\tikzfig{tk_components}
\figcaption{fig:tkcomp}{Resolution is just a right triangle read backwards: the vector is the
hypotenuse, and its components are the two legs. Note which leg gets the cosine; it is always the one
\emph{adjacent} to the marked angle, so the labelling depends on where $\theta$ is measured from.}
\end{figure}

\begin{trapbox}{$\cos$ does not always go with the horizontal}
$A_x=A\cos\theta$ only when $\theta$ is measured \emph{from the $x$-axis}. On a plane inclined at $\theta$
to the horizontal, the weight makes the angle $\theta$ with the \emph{normal} to the slope, not with the
slope itself, so the component along the slope is the one opposite $\theta$: $mg\sin\theta$. Always ask
``adjacent or opposite to the angle I have marked?'' rather than reciting a formula.
\end{trapbox}

\section{Vector Algebra}\label{sec:vecalg}

Adding vectors is the business of \chref{ch:vectors}. What needs flagging now is that vectors (written
$\vc{A}$, with magnitude $A$) can be multiplied in two genuinely different ways, answering two different
questions. Knowing which question a situation asks is more useful than knowing either formula.

\begin{defbox}{The two products, and what each is for}
\textbf{Scalar (dot) product:} returns a \emph{number}:
\[
\vc{A}\cdot\vc{B}=AB\cos\theta=A_xB_x+A_yB_y+A_zB_z,
\]
where $\theta$ is the angle between the vectors drawn tail to tail and $A_x$, $A_y$, $A_z$ are the
components along three perpendicular axes. It measures how much of one vector lies \emph{along} the
other, and for non-zero vectors it vanishes exactly when they are perpendicular. \emph{Used in:} work $W=\vF\cdot\vc{d}$ and power $P=\vF\cdot\vv$ (\chref{ch:energy}).

\medskip
\textbf{Vector (cross) product:} returns a \emph{vector} $\vc{A}\ex\vc{B}$ of magnitude $AB\sin\theta$,
perpendicular to both $\vc{A}$ and $\vc{B}$, its sense fixed by a right-hand rule. Its magnitude is the
\emph{area} of the parallelogram they span, and it vanishes when the vectors are parallel. (The
component forms, the right-hand rule and worked uses of both products are in \S\ref{sec:vecprod}.) \emph{Used in:} torque $\vtau=\vr\ex\vF$ and angular
momentum $\vL=\vr\ex\vp$ (\chref{ch:rotation}).

\medskip
\textbf{The memory hook:} the dot product asks ``how much \emph{along}?'', the cross product ``how much
\emph{across}?''
\end{defbox}

\section{Differentiation}

Most quantities in mechanics are rates: velocity is the rate of change of position, power the rate of doing
work, force the rate of change of momentum. The derivative is the mathematical form of a rate, and only a
handful of its rules are needed.

\begin{defbox}{The derivative is a rate, and a slope}
The derivative of $y$ with respect to $x$ is the limit of a ratio of small changes, where $\Delta x$
(read ``delta $x$'') is a change in $x$ and $\Delta y$ the change in $y$ it produces:
\[
\frac{dy}{dx}=\lim_{\Delta x\to0}\frac{\Delta y}{\Delta x}.
\]
Geometrically it is the \textbf{slope of the tangent} to the graph, the limiting position of a secant
whose two points close in on each other (Figure~\ref{fig:tkderiv}). Physically it is an
\textbf{instantaneous rate}: how fast $y$ changes right now, as opposed to the average change over an
interval. Its units are those of $y$ divided by those of $x$, which is why a rate of change of position
is measured in \si{m/s}. For example, if $x=5t^{2}$ (metres, seconds) then $dx/dt=10t$, so at
$t=\SI{3}{s}$ the instantaneous rate is \SI{30}{m/s}, while the average over the first three seconds is
only $\SI{45}{m}/\SI{3}{s}=\SI{15}{m/s}$. Differentiating twice gives the \textbf{second derivative}
$d^{2}y/dx^{2}$, the rate of change of the rate; acceleration, $d^{2}x/dt^{2}$, is the first example.

The derivatives this book needs:
\[
\frac{d}{dx}\,x^{n}=nx^{n-1},\qquad
\frac{d}{dt}\sin\omega t=\omega\cos\omega t,\qquad
\frac{d}{dt}\cos\omega t=-\omega\sin\omega t,
\]
\[
\frac{d}{dx}e^{kx}=k\,e^{kx},\qquad \frac{d}{dx}\ln x=\frac{1}{x}.
\]
\emph{Used in:} velocity $\vv=d\vr/dt$ and acceleration $\va=d\vv/dt$ (\chref{ch:straightline}), instantaneous
power $P=dW/dt$ (\chref{ch:energy}), and $\vF=d\vp/dt$ (\chref{ch:newton}).
\end{defbox}

\begin{figure}[htbp]
\centering
\tikzfig{tk_derivative}
\figcaption{fig:tkderiv}{Average versus instantaneous. The green \emph{secant} joins two separated points
and gives the average rate $\Delta s/\Delta t$. Sliding the second point inward until it merges with the
first turns the secant into the red \emph{tangent}, whose slope is $ds/dt$, the instantaneous rate.
That limiting process is the whole content of the derivative.}
\end{figure}

\begin{methodbox}{The two combination rules we need}
\textbf{Product rule:}
\[
\frac{d}{dx}(uv)=u\frac{dv}{dx}+v\frac{du}{dx}
\]
\textbf{Chain rule} (for a function of a function: differentiate the outside, then multiply by the
derivative of the inside; a prime denotes a derivative, $f'=df/du$):
\[
\frac{d}{dx}f\big(u(x)\big)=f'\big(u(x)\big)\cdot u'(x)
\qquad\text{e.g.}\qquad
\frac{d}{dt}\sin(\omega t)=\omega\cos(\omega t).
\]
\emph{Used in:} the product rule for differentiating $\vv=v\,\uhat{t}$, speed times a unit vector, which
splits acceleration into its tangential and centripetal parts (\chref{ch:circkin}); the chain rule for
every sinusoid whose argument is not just $t$, and for $a=dv/dt=v\,dv/dx$ (\chref{ch:straightline}).
\end{methodbox}

\begin{defbox}{Maxima and minima: where the slope is zero}
At the top of a hill or the bottom of a valley the tangent is horizontal. So a smoothly varying
quantity $y(x)$ is largest or smallest where
\[
\frac{dy}{dx}=0.
\]
To tell which, look at how the slope changes: if $d^{2}y/dx^{2}<0$ the slope is falling from positive to
negative, and the point is a \textbf{maximum}; if $d^{2}y/dx^{2}>0$ it is a \textbf{minimum}. For
$y=4x-x^{2}$, $dy/dx=4-2x$ vanishes at $x=2$, and $d^{2}y/dx^{2}=-2<0$, so the greatest value is
$y(2)=4$. Two cautions: a zero slope can also be a mere pause ($y=x^{3}$ at $x=0$), and when $x$ is
confined to a range the largest value may sit at an end of it, where the slope need not vanish.
\emph{Used in:} the greatest velocity under a varying acceleration (\chref{ch:straightline}), the least
drift in crossing a river (\chref{ch:relative}), the least force to drag a block (\chref{ch:friction}),
and stable and unstable equilibrium (\chref{ch:energy}).
\end{defbox}

\begin{methodbox}{Logarithmic differentiation: where the error rules come from}
Taking a logarithm turns products into sums, and \emph{then} differentiating gives a rule for
\emph{relative} changes. If
\[
Z=\frac{A^{p}B^{q}}{C^{r}},\qquad\text{then}\qquad \ln Z=p\ln A+q\ln B-r\ln C,
\]
and differentiating each term (using $d(\ln x)=dx/x$, where $dA$ means a very small change in $A$)
gives at once
\[
\frac{dZ}{Z}=p\frac{dA}{A}+q\frac{dB}{B}-r\frac{dC}{C}.
\]
For small errors the differentials can be replaced by the errors themselves, $dA\to\Delta A$ and so on.
The errors in $A$, $B$ and $C$ may have either sign, so in the worst case all three terms push the same
way and their magnitudes add (with $p,q,r$ positive):
\[
\boxed{\;\frac{\Delta Z}{Z}=p\frac{\Delta A}{A}+q\frac{\Delta B}{B}+r\frac{\Delta C}{C}\;}
\]
\emph{Used in:} the combination of errors (\chref{ch:measurement}); this single line \emph{is} the derivation of the
rule that percentage errors add, with powers as multipliers.
\end{methodbox}

\section{Integration}

Integration is used in this book almost entirely in one way: to add up the contributions of many small
pieces of a body, each too small to matter alone. The picture to hold is a sum of thin slices.

\begin{defbox}{An integral is a sum, and an area}
Integration undoes differentiation, but the picture that matters here is different: an integral is the
limit of a \textbf{sum of infinitely many thin slices}. Cutting a region into strips of width $\Delta x$
and height $f(x)$ and adding their areas gives $\sum f(x)\Delta x$; letting the strips become
infinitesimally thin (Figure~\ref{fig:tkint}) turns the sum between $x=a$ and $x=b$ into the \textbf{definite integral}
\[
\int_a^b f(x)\,dx=\lim_{\Delta x\to0}\sum f(x)\,\Delta x,
\]
the area under the graph. The link to differentiation, quoted here from mathematics, is what makes such
sums computable: if $G(x)$ is any function whose derivative is $f(x)$, then
$\int_a^b f(x)\,dx=G(b)-G(a)$. Without limits, the \textbf{indefinite integral} $\int f\,dx=G(x)+C$
carries an arbitrary constant $C$ (not the $C$ of the error rule above), since adding a constant to $G$ leaves its derivative unchanged; in
physics $C$ is fixed by a known starting value, such as the position at $t=0$. For example, $\int_0^{\pi}\sin\theta\,d\theta=[-\cos\theta]_0^{\pi}=2$.
The integrals this book needs:
\[
\int x^{n}\,dx=\frac{x^{n+1}}{n+1}\ (n\neq-1),\qquad
\int_0^{\pi}\sin\theta\,d\theta=2,\qquad
\int_0^{\pi}\sin^{3}\theta\,d\theta=\frac{4}{3}.
\]
\emph{Used in:} work by a variable force $W=\int F\,dx$ (\chref{ch:energy}), and every centre-of-mass and
moment-of-inertia calculation, where the ``slices'' are literal pieces of the body
(\chsref{ch:com}{ch:rotation}).
\end{defbox}

\begin{figure}[htbp]
\centering
\tikzfig{tk_integral}
\figcaption{fig:tkint}{Why an integral is an area. Coarse strips (left) already approximate the region;
making them finer (right) makes the approximation exact in the limit. This is the picture to hold when a
derivation says ``take an element $dm$ and integrate''; you are adding up slices.}
\end{figure}

\begin{methodbox}{The pattern behind every $\int r^2\,dm$}
Whenever this book integrates over a body, three steps recur:
\begin{enumerate}
\item \textbf{Choose a slice} whose every point shares the same value of the quantity being weighted
      (the same distance $r$ from an axis, or the same height $y$).
\item \textbf{Express its mass} $dm$ in terms of a density and the slice's size: $\lambda\,dx$ for a
      wire ($\lambda$ the mass per unit length), $\sigma\,dA$ for a lamina ($\sigma$ the mass per unit
      area, $dA$ here the slice's area), $\rho\,dV$ for a solid ($\rho$ the mass per unit volume).
\item \textbf{Integrate} over the range that sweeps the slice through the whole body.
\end{enumerate}
Choosing the slice well is the entire art; the integration itself is then routine.
\end{methodbox}

\section{Approximations and Algebra}

The last group of tools turns equations into numbers. None is deep, but each comes with a habit of
interpretation that matters as much as the algebra.

\begin{defbox}{The small-angle approximation}
For a small angle measured \textbf{in radians},
\[
\boxed{\;\sin\theta\approx\theta,\qquad \tan\theta\approx\theta,\qquad \cos\theta\approx1-\tfrac12\theta^{2}\approx1\;}
\]
The reason is visible geometrically: for a thin wedge the arc $R\theta$, the opposite side $R\sin\theta$
and the tangent $R\tan\theta$ all shrink towards one another (Figure~\ref{fig:tksmall}, left). The error in $\sin\theta\approx\theta$
stays below one percent up to about $14^\circ$, and in $\tan\theta\approx\theta$ up to about $10^\circ$;
the cruder $\cos\theta\approx1$ is already one percent out at $8^\circ$.
\emph{Used in:} the derivative of a rotating unit vector (\chref{ch:circkin}), the small-angle limit of
the conical pendulum (\chref{ch:circdyn}), and anywhere a curved arc is replaced by a straight chord.
\end{defbox}

\begin{figure}[htbp]
\centering
\tikzfig{tk_smallangle}
\figcaption{fig:tksmall}{Left: in a wedge of angle $\theta$ the perpendicular $R\sin\theta$, the arc $R\theta$ and the tangent $R\tan\theta$ always come in that order, and they close up on one another as the wedge narrows. Right: the three graphs leave the origin with the same slope, which is exactly why the approximation works. In the shaded band ($\theta<10^\circ$) all three agree within one percent; beyond it they part, $\tan\theta$ fastest.}
\end{figure}

\begin{defbox}{The binomial approximation}
For any power $n$ (whole, negative or fractional) and a small number $x$, $|x|\ll1$,
\[
\boxed{\;(1+x)^{n}\approx1+nx\;}
\]
For a whole number $n$ this is the start of the binomial expansion,
$(1+x)^{n}=1+nx+\tfrac12n(n-1)x^{2}+\cdots$, in which every later term carries a further factor of the
small $x$ and can be dropped; that the same first two terms hold for any $n$ is quoted from mathematics.
Geometrically it replaces the curve $(1+x)^{n}$ by its tangent at $x=0$, whose slope is $n$. For example,
$\sqrt{1.02}=(1+0.02)^{1/2}\approx1.01$ (true value $1.00995$), and $1/1.01^{2}\approx1-0.02=0.98$ (true
value $0.9803$). It is the same idea as the error rule above: $(1+\Delta A/A)^{p}\approx1+p\,\Delta A/A$.
\emph{Used in:} the fall of $g$ with height, $(1+h/R)^{-2}\approx1-2h/R$ (\chref{ch:gravitation}),
valid only while $h\ll R$.
\end{defbox}

\begin{keybox}{The quadratic formula}
\[
ax^{2}+bx+c=0\qquad\Longrightarrow\qquad x=\frac{-b\pm\sqrt{b^{2}-4ac}}{2a}
\]
When a physics problem yields two roots, \textbf{both must be examined}: often both are meaningful
(the launch and the landing), and sometimes one is negative or unphysical and is rejected for a stated
reason. Deciding which to keep is part of the physics. \emph{Used in:} finding the time of flight from
$y=ut-\tfrac12gt^{2}$ (\chref{ch:projectiles}), where one root is the launch instant and the other the landing.
\end{keybox}

\begin{methodbox}{Simultaneous equations, and eliminating the unknown you do not want}
Connected-body problems produce one equation per body sharing a common unknown. The usual technique:
\textbf{add or subtract the equations so that the unwanted unknown cancels.} For the Atwood machine
(two masses hung over a pulley, $m_1>m_2$, joined by a string of tension $T$),
\[
m_1g-T=m_1a,\qquad T-m_2g=m_2a,
\]
and \emph{adding} them eliminates $T$ in one step, leaving $(m_1-m_2)g=(m_1+m_2)a$. \emph{Used in:} pulleys
(\chref{ch:connected}) and elastic collisions (\chref{ch:collisions}), where dividing the energy equation by the momentum equation
eliminates the messy squares.
\end{methodbox}

\begin{methodbox}{Proportional reasoning: the fastest tool of all}
Many exam questions ask only \emph{how a quantity scales}, not its value. If $y\propto x^{n}$ then
multiplying $x$ by $k$ multiplies $y$ by $k^{n}$:
\begin{itemize}
\item $P\propto v^{3}$ (the power to push through air whose drag force grows as $v^2$, since power is
      force times speed): double the speed, \emph{eight times} the power.
\item $K=\tfrac12mv^{2}$: double the speed, \emph{four times} the kinetic energy, and, for the same
      braking force, four times the braking distance.
\item $I\propto MR^{2}$ at fixed mass: double the radius, \emph{four times} the moment of inertia.
\end{itemize}
No numbers are needed, only the exponent. \emph{Used in:} drag and power (\chref{ch:energy}), stopping
distances (\chref{ch:straightline}), moments of inertia (\chref{ch:rotation}), and Kepler's third law
(\chref{ch:kepler}).
\end{methodbox}

\section{Summary}\label{sec:summary0}
The chapter assembled twelve tools, and they divide cleanly by purpose. Two \emph{describe}
directions: radian measure, the only angle unit that works inside calculus, and the trigonometric
ratios, used above all to resolve a vector. Two \emph{combine} vectors: the dot product (how much
along) and the cross product (how much across). Three connect a quantity with its rate of change: the derivative extracts an
instantaneous rate from a changing quantity (its product, chain and logarithmic forms extend it), its
zeros locate maxima and minima, and the integral adds up infinitely many thin slices (one slicing
pattern serves every body). The last five turn equations into answers: the small-angle and binomial
approximations, the quadratic formula (where both roots must be
examined for physical meaning), the elimination of an unwanted unknown between simultaneous equations,
and proportional reasoning.

\begin{itemize}
\item Radians are dimensionless and $s=R\theta$; an angle standing on its own in a formula must be in
      radians.
\item $A_x=A\cos\theta$ only when $\theta$ is measured from the $x$-axis; cosine goes with the
      \emph{adjacent} leg.
\item $\vc{A}\cdot\vc{B}=AB\cos\theta$ is a number; $|\vc{A}\ex\vc{B}|=AB\sin\theta$ is the area of the
      parallelogram the two vectors span.
\item Logarithmic differentiation turns $Z=A^pB^q/C^r$ into $dZ/Z=p\,dA/A+q\,dB/B-r\,dC/C$, whose
      worst case, $\Delta Z/Z=p\,\Delta A/A+q\,\Delta B/B+r\,\Delta C/C$, is the whole theory of error
      propagation in one line.
\item A smooth quantity is greatest or least where its derivative is zero; the sign of the second
      derivative says which. For small $x$, $(1+x)^n\approx1+nx$.
\item Every integral over a body follows one pattern: choose a slice on which the weighted quantity is
      constant, express its mass, and integrate.
\item Proportional reasoning ($y\propto x^n$, so scaling $x$ by $k$ scales $y$ by $k^n$) answers many
      examination questions with no arithmetic at all.
\end{itemize}

\section{Problem Types}\label{sec:types0}
Mathematics is rarely examined for its own sake, but four slips recur so often inside physics questions
that they are worth isolating: an angle left in degrees, a missing chain-rule factor, a variable
quantity treated as constant, and a confusion between a factor and an increase. Each problem type below
gives the \emph{tell} that identifies the situation, the \emph{method}, and the \emph{trap}.

\archq{Arc length and the radian}
{An angle is given in degrees, and the formula uses it on its own: $s=R\theta$, $v=\omega r$,
$\omega=\theta/t$.}
{Convert first: $\theta_{\text{rad}}=\theta^\circ\times\pi/180$. Then substitute.}
{Substituting the number of degrees directly, which gives an answer about $57$ times too large.}
{A wheel of radius \SI{0.50}{m} turns through $60^\circ$. The distance moved by a point on its rim is:}
{\SI{30}{m}}{\SI{0.52}{m}}{\SI{0.26}{m}}{\SI{1.05}{m}\archans{B}{$s=R\theta=(0.50)(\pi/3)=\SI{0.52}{m}$.
Option (A) uses $\theta=60$ directly; option (D) uses the diameter in place of the radius.}}

\archq{Differentiating to find a rate}
{A quantity is given as a function of time, such as $x(t)$, and its rate at a particular instant is
wanted.}
{Differentiate term by term; for $\sin(\omega t)$ or $\cos(\omega t)$ the chain rule brings out a factor
$\omega$. Then substitute the instant.}
{Dividing $x$ by $t$, which gives an \emph{average} rate, or forgetting the factor $\omega$.}
{A particle moves along a line with $x=2\sin(3t)$, where $x$ is in metres and $t$ in seconds. Its
velocity at $t=0$ is:}
{$0$}{\SI{2}{m/s}}{\SI{6}{m/s}}{\SI{3}{m/s}\archans{C}{$v=dx/dt=6\cos(3t)$, which is \SI{6}{m/s} at
$t=0$. Option (A) is the position at $t=0$; option (B) omits the chain-rule factor $3$.}}

\archq{Integrating a varying quantity}
{The quantity to be summed changes along the way: the area under a curve, or the total of something
whose rate is a given function.}
{Write the sum as a definite integral and evaluate it as $G(b)-G(a)$, where $G'=f$.}
{Multiplying the final value of $f$ by the interval, as though $f$ had been constant throughout.}
{The area under the graph of $y=6x^2$ between $x=0$ and $x=2$ is:}
{$8$}{$16$}{$24$}{$48$\archans{B}{$\int_0^2 6x^2\,dx=\big[2x^3\big]_0^2=16$. Option (D) is
$y(2)\times2=24\times2$, the constant-height trap; option (C) is $y(2)$ itself.}}

\archq{Scaling by proportional reasoning}
{``By what factor (or what percentage) does $y$ change if $x$ changes by \ldots?'', with a power law
relating them.}
{If $y\propto x^n$, scaling $x$ by $k$ scales $y$ by $k^n$. Convert the answer to a percentage
\emph{increase} only at the end: an increase is $(k^n-1)\times100\%$.}
{Reporting the factor $k^n$ as though it were the percentage increase, or using $nk$ in place of $k^n$.}
{The speed of a body of fixed mass is increased by $50\%$. Its kinetic energy increases by:}
{$50\%$}{$100\%$}{$125\%$}{$225\%$\archans{C}{$K\propto v^2$ and $k=1.5$, so $K$ is multiplied by
$1.5^2=2.25$, an increase of $125\%$. Option (D) mistakes the factor for the increase; option (B) is
$n\times50\%$.}}
